\section{Chain of Responsibility}
\label{sec:pat-chain}

\problemstatement{Pass a request along a chain of handlers; each handler
decides either to process the request or to forward it to the next,
decoupling the sender from the specific receiver that ends up handling it.}

\begin{patternbox}[When You Reach for It]
The trigger is \emph{``a request that several handlers might process, and
we don't know in advance which one will.''} Approval workflows, event
bubbling, logging levels, and middleware pipelines all fit: try each
handler in turn until one takes it.
\end{patternbox}

\subsection*{Understanding the pattern}

Each handler holds a pointer to the next and implements a
\texttt{handle(request)}: if it can deal with the request it does;
otherwise it delegates to the next in line. The client fires the request
at the head of the chain and never needs to know which link resolves it.
Handlers can be reordered or inserted without touching the sender.

\begin{ideabox}[Handle or Pass Along]
Give every handler the same interface and a ``next'' link. The decision
``mine or not mine'' is local to each handler; the chain's shape is
configured externally. This decouples \emph{who sends} from \emph{who
handles}.
\end{ideabox}

\subsection*{C++ implementation}

\begin{lstlisting}
#include <bits/stdc++.h>
using namespace std;

class Approver {
protected:
    Approver* next = nullptr;
public:
    Approver* setNext(Approver* n) { next = n; return n; }
    virtual void handle(int amount) {
        if (next) next->handle(amount);        // default: pass along
        else cout << "No one could approve " << amount << "\n";
    }
    virtual ~Approver() = default;
};

class Manager : public Approver {
public:
    void handle(int amount) override {
        if (amount <= 1000) cout << "Manager approves " << amount << "\n";
        else Approver::handle(amount);
    }
};
class Director : public Approver {
public:
    void handle(int amount) override {
        if (amount <= 10000) cout << "Director approves " << amount << "\n";
        else Approver::handle(amount);
    }
};

int main() {
    Manager mgr; Director dir;
    mgr.setNext(&dir);                          // build the chain
    mgr.handle(500);                            // Manager
    mgr.handle(5000);                           // passed to Director
    mgr.handle(50000);                          // unhandled
    return 0;
}
\end{lstlisting}

\begin{complexitybox}[Trade-offs]
\textbf{Pros.} Decouples sender and receiver; handlers are added,
removed, or reordered freely; each handler is simple. \textbf{Cons.} A
request may fall off the end unhandled; harder to trace which handler
acted; a long chain adds latency. Ensure a terminal/default handler.
\end{complexitybox}

\begin{takeawaybox}
Chain of Responsibility lets a request travel handler-to-handler until one
processes it, decoupling the caller from the resolver. Recognise it in
approval hierarchies and middleware. The key risk to mention: unhandled
requests, so design an end-of-chain default.
\end{takeawaybox}
